\documentclass{article}
\usepackage{amsmath, amssymb, amsthm}
\usepackage{hyperref}
\usepackage{geometry}
\geometry{margin=1in}

\title{Erd\H{o}s Problem \#3}
\date{Last updated: 2025-08-31}
\author{Source: erdosproblems.com}

\begin{document}
\maketitle

\noindent\textbf{Status:} OPEN \quad \textbf{Prize: \$5000}

\noindent\textbf{Tags:} number theory, additive combinatorics, arithmetic progressions

\medskip

\noindent\textbf{Problem Statement:}

If $A\subseteq \mathbb{N}$ has $\sum_{n\in A}\frac{1}{n}=\infty$ then must $A$ contain arbitrarily long arithmetic progressions?




This is essentially asking for good bounds on $r_k(N)$, the size of the largest subset of $\{1,\ldots,N\}$ without a non-trivial $k$-term arithmetic progression. For example, a bound like\[r_k(N) \ll_k \frac{N}{(\log N)(\log\log N)^2}\]would be sufficient. Even the case $k=3$ is non-trivial, but was proved by Bloom and Sisask \cite{BlSi20}. Much better bounds for $r_3(N)$ were subsequently proved by Kelley and Meka \cite{KeMe23}. Green and Tao \cite{GrTa17} proved $r_4(N)\ll N/(\log N)^{c}$ for some small constant $c&gt;0$. Gowers \cite{Go01} proved\[r_k(N) \ll \frac{N}{(\log\log N)^{c_k}},\]where $c_k&gt;0$ is a small constant depending on $k$. The current best bounds for general $k$ are due to Leng, Sah, and Sawhney \cite{LSS24}, who show that\[r_k(N) \ll \frac{N}{\exp((\log\log N)^{c_k})}\]for some constant $c_k&gt;0$ depending on $k$. Curiously, Erd\H{o}s \cite{Er83c} thought this conjecture was the 'only way to approach' the conjecture that there are arbitrarily long arithmetic progressions of prime numbers, now a theorem due to Green and Tao \cite{GrTa08} (see [219] ). In \cite{Er81} Erd\H{o}s makes the stronger conjecture that\[r_k(N) \ll_C\frac{N}{(\log N)^C}\]for every $C&gt;0$ (now known for $k=3$ due to Kelley and Meka \cite{KeMe23}) - see [140] . See also [139] and [142] . This is discussed in problem A5 of Guy's collection \cite{Gu04}.




References

[BlSi20] Bloom, T.F. and Sisask, O., Breaking the logarithmic barrier in Roth's theorem on arithmetic progressions . arXiv:2007.03528 (2020).

[Er81] Erd\H{o}s, P., On the combinatorial problems which I would most like to see solved . Combinatorica (1981), 25-42.

[Er83c] Erd\H{o}s, Paul, Combinatorial problems in geometry . Math. Chronicle (1983), 35-54.

[Go01] Gowers, W. T., A new proof of Szemer\'{e}di's theorem . Geom. Funct. Anal. (2001), 465-588.

[GrTa08] Green, Ben and Tao, Terence, The primes contain arbitrarily long arithmetic progressions . Ann. of Math. (2) (2008), 481-547.

[GrTa17] Green, Ben and Tao, Terence, New bounds for Szemer\'{e}di's theorem, III: a polylogarithmic bound for $r_4(N)$ . Mathematika (2017), 944-1040.

[Gu04] Guy, Richard K., Unsolved problems in number theory . (2004), xviii+437.

[KeMe23] Kelley, Z. and Meka, R., Strong Bounds for 3-Progressions . arXiv:2302.05537 (2023).

[LSS24] Leng, J., Sah, A. and Sawhney, M., Improved bounds for Szemer\'{e}di's theorem . arXiv:2402.17995 (2024).

\noindent\textbf{Related OEIS sequences:} A003002, A003003, A003004, A003005


\bigskip
\noindent\small{Source: \url{https://www.erdosproblems.com/3}}
\end{document}
